%---------------------------------------------------------------
\section{Postup při realizaci projektu}
%---------------------------------------------------------------

% Outline:
% - Výchozí stav projektu a důvody modernizace.
% - Knihovny @eduriam/ui-core a @eduriam/ui-x.
% - Aktualizace technologického stacku.
% - Refaktoring kódu.
% - Aktualizace linterů a konvencí.
% - Zavedení AI konvencí.
% - Nový Figma prototyp.

\subsection{Příprava realizace projektu}

Před zahájením implementace jednotlivých funkcí jsem navázal na původní frontendové řešení z bakalářské práce a provedl jsem jeho cílenou modernizaci. Důvodem nebyla pouze potřeba doplnit nové funkce, ale i zásadní změna domény aplikace a požadavek na dlouhodobě udržitelnou strukturu kódu. Hlavním cílem přípravy tedy bylo oddělit stabilní stavební části systému, aktualizovat technologický základ a odstranit omezení, která komplikovala další rozvoj.

\subsubsection{Příprava knihoven \texttt{@eduriam/ui-core} a \texttt{@eduriam/ui-x}}

V rámci přípravy jsem oddělil komponentovou vrstvu do dvou knihoven. Knihovna \texttt{@eduriam/ui-core} slouží jako základ design systému postavený nad Material UI\footnote{Dostupné na: \href{https://mui.com/}{https://mui.com/}}. Obsahuje zejména základní stavební prvky uživatelského rozhraní a centrální definici \textit{theme}, která určuje konzistentní vzhled aplikace.

Knihovna \texttt{@eduriam/ui-x} obsahuje komplexnější komponenty, které se využívají zejména při učení a vysvětlování látky. Tyto komponenty jsou navržené tak, aby byly znovupoužitelné i mimo hlavní frontend aplikace. Praktickým přínosem je možnost samostatně zobrazit a testovat průběh učení bez nutnosti vázat celý obsah na konkrétní implementaci jedné stránky.

Ve frontendu jsem následně používal obě knihovny současně. Díky tomu jsem oddělil design systém od aplikační logiky a vytvořil předpoklad pro konzistentní využití komponent i v dalších částech projektu.

\subsubsection{Aktualizace knihoven a nástrojů}

Technologický základ jsem aktualizoval na novější verze, konkrétně například Yarn 4.9.3, Node.js 24.9.0, Storybook 10.3.1, Next.js 16.2.1 a Material UI 7.3.9. Tyto změny byly důležité zejména kvůli kompatibilitě mezi knihovnami, dostupnosti aktuálních funkcí a snížení technického dluhu v projektu.

Aktualizace zároveň zjednodušila další implementační kroky, protože jsem pracoval nad konzistentním prostředím, které odpovídalo současným doporučeným postupům.

\subsubsection{Refaktoring kódu}

Po aktualizaci stacku jsem provedl refaktoring stávajícího kódu. Nejprve jsem odstranil části navázané na původní doménu výuky jazyků, které již neodpovídaly novému zaměření aplikace.

Dále jsem přepracoval strukturu komponent pro učení tak, aby bylo možné jednoduše přidávat nové typy cvičení a vysvětlujících bloků. Původní řešení bylo příliš rigidní a rozšiřování funkcí vedlo k neúměrnému nárůstu složitosti. Nová struktura je modulárnější a lépe podporuje iterativní rozvoj.

Současně jsem zavedl sjednocující layout komponenty, zejména \texttt{PageRoot}, \texttt{ContentContainer} a \texttt{PageNavigation}. Tyto prvky zajišťují konzistentní rozložení uživatelského rozhraní napříč stránkami, například jednotné odsazení, responzivní chování a jednotný způsob definice navigačních prvků.

\subsubsection{Aktualizace linterů a konvencí}

V projektu jsem aktualizoval ESLint a Prettier a upravil jejich konfiguraci tak, aby byla kontrola kvality kódu přísnější a zároveň čitelná pro běžný vývoj. Praktickým příkladem je zavedení pravidla \texttt{eqeqeq}\footnote{Dostupné na: \href{https://eslint.org/docs/latest/rules/eqeqeq}{https://eslint.org/docs/latest/rules/eqeqeq}}, které vynucuje porovnávání pomocí \texttt{===} a \texttt{!==}. Tím se snižuje riziko nejednoznačného porovnávání hodnot a chyb způsobených implicitní konverzí typů v jazyce JavaScript/TypeScript.

\subsubsection{Zavedení nástrojů umělé inteligence}

Další částí přípravy bylo zavedení konvencí pro AI asistenty používané při vývoji. Vycházel jsem ze standardu Agent Skills\footnote{Dostupné na: \href{https://agentskills.io/}{https://agentskills.io/}}. V rámci projektu jsem definoval pravidla zejména pro generování nových komponent, stránek a \texttt{.feature} souborů v jazyce Gherkin.

Výhodou tohoto přístupu je, že stejné konvence lze využívat napříč různými nástroji a modely, například v prostředích Claude Code, Cursor, Codex nebo Gemini CLI. Díky tomu jsem snížil rozdíly v kvalitě výstupu při použití různých AI nástrojů.

\subsubsection{Vytvoření nového Figma prototypu}

Původní prototyp byl založen na externí lo-fi knihovně komponent. Tento přístup byl vhodný pro rychlé vytvoření prvotního návrhu, ale v pokročilé fázi návrhu komplikoval práci, protože struktura komponent a jejich parametry neodpovídaly cílové implementaci.

Z tohoto důvodu jsem vytvořil nový prototyp s ručně navrženými komponentami bez závislosti na externí knihovně. Výsledkem byla nižší komplexita návrhu, přesnější návaznost na implementaci a jednodušší převod návrhu do produkčního kódu.
